\section{Challenges and Solutions}
\label{sec:challenges}

\noindent\textbf{Reliability and Accuracy}

\textbf{Challenge:} Ensuring reliability and accuracy in a natural-language Slurm agent is difficult because Slurm commands depend on exact syntax, object identifiers, flags, partitions, accounts, and current scheduler state. A generated answer may sound plausible while still selecting the wrong tool, referring to the wrong job, misunderstanding a pending reason, or confusing documentation knowledge with live cluster facts.

\textbf{Solutions:} The system grounds user requests in typed MCP tools instead of free-form shell generation. The agent uses a ReAct-style execution pattern \cite{yao2022react}, but the important project-specific layer is the Slurm grounding around it: live state is retrieved through Slurm tools, command and state explanations are supported by local official-documentation retrieval, and under-specified requests can trigger clarification rather than forced execution.

\noindent\textbf{Safety in Cluster Operations}

\textbf{Challenge:} Slurm is a shared operational environment. Read-only inspection is low risk, but cancelling jobs, holding jobs, changing node state, modifying reservations, or updating accounting state can interrupt users and affect shared resources. These actions cannot be handled like ordinary chatbot responses.

\textbf{Solutions:} The system separates read-only observation from state-changing operation through an Observer/Operator workflow. The Observer handles inspection, diagnosis, accounting/resource-limit queries, reservation checks, documentation retrieval, and web-search fallback. The Operator handles state-changing work through structured handoffs and a pending-action workflow that requires explicit user confirmation before execution. Dynamic tool filtering keeps sub-agents within their intended responsibility boundaries, and MCP parameter validation prevents arbitrary shell-command construction.

\noindent\textbf{Evaluation and Dataset Construction}

\textbf{Challenge:} Evaluating an agentic Slurm assistant requires more than text-quality scoring. There is no ready-made public dataset covering the exact Slurm workflows and tool behavior needed; the benchmark must be manually designed as synthetic scenarios spanning read-only monitoring, diagnosis, submission, job control, accounting, documentation questions, multi-step workflows, and edge cases. Beyond benchmark construction, reliable automated evaluation additionally requires: (a)~deterministic source-state resets between test cases — without them, state contamination makes pass/fail scores non-reproducible; and (b)~structural traces of tool calls, routing decisions, and HITL activations — because a response can describe the correct operation while still calling the wrong tool, skipping a required handoff, or failing to trigger confirmation.

\textbf{Solutions:} The project manually constructs a 3,135-case scenario-grid benchmark across eleven balanced categories. Each case specifies expected tool behavior, routing, confirmation policy, and source-to-target state where applicable. Mock Slurm scenarios can be reset before each case; parallel evaluation workers use isolated session identifiers and worker-specific MCP endpoints. The evaluation harness records full tool and routing traces, so failures are diagnosable rather than reduced to an opaque aggregate score.

\noindent\textbf{Diversity of Slurm Workflows}

\textbf{Challenge:} Slurm administration spans many domains: jobs, nodes, partitions, accounting, QoS, fair-share, reservations, configuration, scheduler health, job arrays, dependencies, and selected administrative actions. A narrow prototype covering only \texttt{squeue} and \texttt{scancel} would not represent the operational breadth expected from an HPC assistant.

\textbf{Solutions:} The MCP tool layer exposes a broad Slurm operation surface, and the benchmark mirrors that breadth across eleven balanced categories. Live or mock MCP command results remain the source of truth for cluster-state questions, while local documentation retrieval handles reference-heavy questions, keeping live facts separate from model memory.

\noindent\textbf{Domain-Adapted Fine-Tuning}

\textbf{Challenge:} Relying on commercial API models introduces per-request cost, latency, and data-privacy concerns for production HPC environments. However, smaller open-weight models suffer from two compounding limitations: (1)~they do not natively understand Slurm tool-calling conventions, and (2)~their tool-selection accuracy degrades as the number of available tools grows---a phenomenon documented in recent tool-use scaling studies~\cite{qin2023toolllm}. The Slurm agent's combined tool surface comprises 64 agent-facing tools when presented to a single monolithic agent; this leads to selection confusion, where the model calls \texttt{sacct} when \texttt{squeue} is appropriate or hallucinates non-existent commands such as \texttt{scontrol\_requeue}.

\textbf{Solutions:} The project addresses both limitations through a dual strategy. First, the Observer/Operator architecture partitions the tool surface into two focused subsets---29 read-only tools for the Observer and 40 tools for the Operator (35 action tools plus 5 discovery reads)---reducing each agent's selection space from 64 and recovering tool-recall accuracy (empirically +10.2 percentage points on routing-neutral tasks, Section~\ref{sec:ablation}). Second, the project distills successful evaluation traces from the commercial model into a structured training corpus, then applies parameter-efficient fine-tuning---LoRA adapters on a 4-bit quantized base (QLoRA)---to an open-weight model (Qwen2.5-14B). Role-scoped tool exposure ensures each training sample includes only the tools available to its role, so the model internalizes both the safety boundary and the reduced tool surface as part of its generation behavior. The resulting model runs entirely on local GPU infrastructure with no external API dependency.
